\documentclass[10pt,a4paper]{amsart}
\usepackage[margin=19mm]{geometry}
\usepackage{amsmath,amssymb,mathtools}
\usepackage[scr=rsfs,cal=euler]{mathalfa}
\usepackage{xfrac}
\usepackage[unicode,hidelinks,bookmarks=false]{hyperref}
\newcommand{\ie}{i.e.\ }
\newcommand{\cf}{cf.\ }
\newcommand{\resp}{respectively\ }
\newcommand{\Subsection}{\subsection}
\providecommand{\nobreakdash}{\nobreak\hbox{-}}
\setlength{\emergencystretch}{2em}
\multlinegap0pt
\usepackage{xcolor}
\usepackage{amssymb}
\usepackage[shortlabels]{enumitem}
\newtheorem{lem}{Lemma}
\def\N{{I\!\!N}}
\title{Multiplicity of solutions to a class of degenerate elliptic equations in both sub-critical and critical cases}
\author{Kaushik Bal, Sanjit Biswas}
\date{Source: arXiv:2412.04794v2 (CC BY 4.0)}
\begin{document}
\maketitle

\noindent\emph{Source and licence.} Multiplicity of solutions to a class of degenerate elliptic equations in both sub-critical and critical cases. Version arXiv:2412.04794v2 (CC BY 4.0), distributed by arXiv under the Creative Commons Attribution 4.0 International licence, http://creativecommons.org/licenses/by/4.0/, as recorded in the arXiv licence metadata for this exact version and shown on the abstract page https://arxiv.org/abs/2412.04794v2. Full licence notice: \texttt{licenses/arxiv-2412.04794-CC-BY-4.0.txt}.\par\smallskip

\noindent\textit{Editorial scope.} Counted unit: the article's lem BrezisNiren1 (source0.tex lines 677--679). The article's complete original proof of it is delivered with it (source0.tex lines 681--744, one \texttt{proof} environment of 64 lines). No mathematics is written, repaired or re-derived: the statement and the proof are the article's own lines, in the article's own notation, and no retained line is substituted or re-flowed.  Retained uncounted as the source-local context the proof needs: lines 161--166; lines 624--626.  Nothing was omitted from any retained span, so the package carries no omission record.  No retained line contains a citation, so the wrapper carries no bibliography and no bibliography entry.  No retained line contains a figure environment, an image-inclusion command or drawing code, so no image is delivered and the package declares no asset dependency.  Editorial formatting only, in addition: an AMS \texttt{amsart} preamble that restates the article's own declarations and macros used by the retained lines, namely the environments \texttt{lem} on separate counters, together with the standard packages those lines need.  The retained environments print in this document as Lemma 1 in order of appearance, whereas the article prints the same items with the numbering of its own full document.  The complete native source of the article as distributed in the arXiv e-print for this exact version is bundled unchanged as \texttt{sources/169702/source0.tex}.
\begin{align}\label{ME1}
    \left. \begin{array}{l}
        -\Delta_\lambda u=\mu g|u|^{r-1}u+|u|^{2^*_\lambda-2}u \;\text{in}\; \Omega \\
         u\in H^{1,\lambda}_0(\Omega)
    \end{array}\right\} \tag{$E_\mu$}
\end{align}

\begin{lem}\label{Fsol}
     There exists $\mu^*>0$ such that for every $\mu\in(0,\mu^*)$, the \textcolor{black}{problem (\ref{ME1})} admits at least one non-zero solution $u_\mu$ with $I_\mu(u_\mu)<0$.
\end{lem}

\begin{lem}\label{BrezisNiren1}
     The functional $I_\mu$ satisfies $(PS)_c$ condition for every $c<\Tilde{c}$ and $\mu\in(0,\mu^*)$. Here $\mu^*$ is obtained in Lemma \ref{Fsol}.
\end{lem}

\begin{proof}
Let $c<\Tilde{c}$ and $\{u_n\}$ be a sequence in $H^{1,\lambda}_0(\Omega)$ such that 
\begin{align}\label{W5}
    \begin{cases}
        I_\mu(u_n)\to c \text{ as } n\to\infty\\
        I'_\mu(u_n)\to 0 \text{ in } (H^{1,\lambda}_0(\Omega))^*.
    \end{cases}
\end{align}
Applying this fact, we deduce that for large enough $n\in\N,$
\begin{align*}
    \frac{1}{Q}\|u_n\|^2_\lambda-\mu\left (\frac{1}{r+1}-\frac{1}{2^*_\lambda}\right )\int_\Omega g\cdot(u_n^+)^{r+1}=I_\mu(u_n)-\frac{1}{1+s}\langle I'_\mu(u_n),u_n\rangle\leq c+1+m\|u_n\|_\lambda,
    \end{align*}
  for some $m>0$. Using H\"{o}lder inequality, we obtain
    \begin{align}\label{BDD}
         \frac{1}{Q}\|u_n\|_\lambda^2-\mu\left (\frac{1}{r+1}-\frac{1}{2^*_\lambda}\right )\|g\|_a\|u_n\|_\lambda^{1+r}\leq c+1+m\|u_n\|_\lambda,
    \end{align}
   which yields that the sequence $\{u_n\}$ is bounded in $H^{1,\lambda}_0(\Omega)$. Consequently, we can assume that there exists $u\in H^{1,\lambda}_0(\Omega)$ such that up to a subsequence, the following holds.
   \begin{align}\label{Lab}
        \begin{cases}
       u_n\rightharpoonup u \text{ weakly in }H^{1,\lambda}_0(\Omega),\\
        u_n\to u \text{ strongly in } L^{q}(\Omega) \text{ for all }1\leq q< 2^*_\lambda,\\
        u_n\to u \text{ pointwise a.e. in }\Omega.
    \end{cases}
   \end{align}
 Using this fact, one has
\begin{align}\label{W3}
    o(1)=\langle I'_\mu(u_n),\phi\rangle=\langle I'_\mu(u),\phi\rangle+o(1) \text{ for all }\phi\in H^{1,\lambda}_0(\Omega).
\end{align}
Thus, $u$ is a critical point of $I_\mu$, and hence non-negative. Since $c<\Tilde{c}$, from the definition of $\Tilde{c}$, we have $c<I_\mu(u)+\frac{1}{Q}\mathcal{S}_\lambda^\frac{Q}{2}$. Next we prove that the sequence $\{u_n\}$ has a convergent subsequence in $H^{1,\lambda}_0(\Omega)$. We have
\begin{align}
    &\|u_n\|_\lambda^2=\|u_n-u\|_\lambda^2+\|u\|_\lambda^2+o(1)\label{W1},\\
    \text{ and }&\nonumber\\
    &\langle I_\mu'(u_n),u_n\rangle=\int_\Omega |\nabla_\lambda u_n|^2\, dz-\mu\int_\Omega g\cdot(u_n^+)^{1+r} \, dz-\int_\Omega(u_n^+)^{2^*_\lambda} dz=o(1).\label{TA1}
\end{align}
Thus, it is immediate from the facts (\ref{TA1}) and $\langle I'_\mu(u),u\rangle=0$ that
\begin{align}\label{W2}
    \int_\Omega |\nabla_\lambda u_n|^2\, dz-\int_\Omega|u_n|^{2^*_\lambda} dz=\mu\int_\Omega g(u^+)^{1+r} \, dz + o(1)=\int_\Omega |\nabla_\lambda u|^2\, dz-\int_\Omega(u^+)^{2^*_\lambda} dz+o(1).
\end{align}
Since $\{u_n^+\}$ is bounded in $L^{2^*_\lambda}(\Omega)$, applying Brezis-Lieb Lemma, one has 
\begin{align}\label{W4}
  \|u_n^+\|_{2^*_\lambda}^{2^*_\lambda} = \|(u_n-u)^+\|_{2^*_\lambda}^{2^*_\lambda}+\|u^+\|_{2^*_\lambda}^{2^*_\lambda}+o(1).
\end{align}
The identity (\ref{W1}), (\ref{W2}) and (\ref{W4}) imply that
\begin{align}\label{W8}
    \|u_n-u\|_\lambda^2=\|(u_n-u)^+\|_{2^*_\lambda}^{2^*_\lambda}+o(1)
\end{align}
We substitute (\ref{W1}) and (\ref{W4}) in the expression of $I_\mu(u_n)$ and using the fact (\ref{W5}), we obtain 
\begin{align*}
    \frac{1}{2}\left [\|u\|_\lambda^2+\|u_n-u\|_\lambda^2 \right ]-\frac{\mu}{r+1}\int_\Omega g\cdot(u_n^+)^{1+r}\, dz-\frac{1}{1+s}\left [\|(u_n-u)^+\|_{2^*_\lambda}^{2^*_\lambda}+\|u^+\|_{2^*_\lambda}^{2^*_\lambda}\right ]=c+o(1),
\end{align*}
which implies
\begin{align}\label{W9}
    I_\mu(u)+\left [\frac{1}{2}\|u_n-u\|_\lambda^2-\frac{1}{1+s} \|(u_n-u)^+\|_{2^*_\lambda}^{2^*_\lambda} \right ]=c+o(1).
\end{align}
Suppose $u_n$ does not converge to $u$ in $H^{1,\lambda}_0(\Omega)$, i.e., $\|u_n-u\|_\lambda\to d\neq 0$. From the definition of $\mathcal{S}_\lambda$, we get

\begin{align}
    \|u_n-u\|_\lambda^2\geq \mathcal{S}_\lambda\|u_n-u\|^2_{2^*_\lambda}\geq \mathcal{S}_\lambda\|(u_n-u)^+\|^2_{2^*_\lambda}
\end{align}
This fact, together with the identity (\ref{W8}), ensures that $d$ enjoys the following inequality:
$$d^2\geq \mathcal{S}_\lambda d^\frac{4}{2^*_\lambda}.$$
Thus, $d\geq \mathcal{S}_\lambda^\frac{Q}{4}$. Using this fact, we can infer from the identity (\ref{W9}) that $$c=I_\mu(u)+\frac{1}{Q}d^2\geq I_\mu(u)+\frac{1}{Q}\mathcal{S}_\lambda^\frac{Q}{2}$$
which is a contradiction as $c<I_\mu(u)+\frac{1}{Q}\mathcal{S}_\lambda^\frac{Q}{2}$. Consequently, $u_n\to u$ strongly in $H^{1,\lambda}_0(\Omega)$.
\end{proof}

\end{document}
